\section{Prossimi passi}\label{sec:next}
\begin{enumerate}[label=\textbf{N\arabic*.},leftmargin=2.2em]
\item \textbf{Legge generale dei cumulanti.} Dimostrare sul reticolo $NC(k)$ che il termine $\kappa_\pi(A)\varphi_{K(\pi)}(B)$ decade come $(1-p)^{(k-\#\text{singoletti}(\pi))t}$. Il meccanismo all'ordine dominante in $1/D$ è già nella \cref{sec:bottleneck}.
\item \textbf{Esperimenti.} In un esperimento di eco con rumore regolabile, i tassi $(1-p)^{n}$ separano i cumulanti liberi. L'idea è usare la decoerenza come spettroscopia della struttura libera dell'OTOC; i dati di eco dei processori superconduttori sono il banco di prova naturale.
\item \textbf{Rumore locale nel bagno a $D$ finito.} Collegare i due modelli: un bagno a collo di bottiglia con $L$ qudit rumorosi dovrebbe mostrare l'incrocio del modello 1, controllato da $pL$.
\item \textbf{Tempo continuo.} Un bagno Browniano dà il generatore esatto a $k$ repliche (domanda Q3 delle note), e permette di verificare l'inferenza che $\mathcal K$ sia di ordine $dt$.
\item \textbf{Quenched contro annealed.} Stabilire se gli incroci sopravvivono nelle quantità tipiche (limite di repliche $k\to1$), come nei lavori su RMPS.
\item \textbf{Randomness strutturata e doping.} Ripetere con bagni Clifford random (commutante di Clifford) e Gaussiani fermionici, poi con doping $T$ o non Gaussiano.
\end{enumerate}
